AI-simulated expert panels have been proposed as scalable, low-cost proxies for human expert judgment in creative work evaluation. A critical prerequisite for responsible deployment is calibration: systematic evidence of which quality dimensions AI scores reliably, which it underestimates, and by how much. We present the first calibration study of AI expert panel scores against human expert judgment in a creative domain.

In a controlled study, 10 puzzle hunt puzzles were independently scored by an AI panel (using v1 and v2 expert persona profiles) and five experienced human puzzle hunt constructors on an identical six-dimension rubric (Clarity, Solvability, Elegance, Reading Reward, Fun, Confirmation Quality). We compute per-dimension Pearson and Spearman correlations, measure signed bias (AI mean minus human mean), test whether richer AI profiles close the calibration gap, and derive corrective scaling factors for practitioners.

Key findings: (1) AI panel scores correlate strongly with human judgment on \emph{technical} dimensions---Clarity ($r = 0.88$), Solvability ($r = 0.84$), and Confirmation Quality ($r = 0.86$)---where quality criteria are verifiable and non-hedonic. (2) Correlation is substantially weaker on \emph{hedonic} dimensions---Fun ($r = 0.52$) and Elegance ($r = 0.68$)---where evaluation requires the experience of solving, not merely inspection of structure. (3) AI panels systematically underestimate Fun ($\Delta = -0.83$, $p < 0.01$) and Elegance ($\Delta = -0.58$, $p < 0.05$) relative to human panels; technical dimensions show no significant bias. (4) Richer v2 profiles substantially close the Elegance gap but leave the Fun gap essentially unchanged, confirming that Elegance has a structural component that persona depth can capture, while Fun requires hedonic experience that no profile enrichment can substitute for.

We provide practitioner-ready corrective scaling factors for all six dimensions, with confidence intervals, and derive a recalibrated pass threshold for AI-only review pipelines that accounts for the systematic hedonic underestimation. Our results generalize the principle that AI evaluation reliability tracks criterion verifiability: AI panels are appropriate gatekeepers for technically verifiable quality, and should be supplemented with human review for hedonic dimensions.
